\documentclass[11pt]{article}
\usepackage[margin=1in]{geometry}
\usepackage{amsmath,amssymb}
\usepackage{booktabs}
\usepackage{graphicx}
\usepackage{hyperref}
\usepackage{listings}
\usepackage{xcolor}
\usepackage{float}
\usepackage{array}

\hypersetup{colorlinks=true, linkcolor=blue, urlcolor=blue}

\lstset{
  basicstyle=\ttfamily\small,
  breaklines=true,
  frame=single,
  backgroundcolor=\color{gray!10},
  columns=fullflexible,
}

\title{CAROM: Content-Addressable Recurrent Operator Machine \\
\large GPU Mechanism Screen Report --- Experiment r1}
\author{Research Agent}
\date{2026-07-20}

\begin{document}
\maketitle

\begin{abstract}
CAROM is a differentiable register machine that separates control (command tokens route computation via a learned routing table) from data (workspace slots operated on by a shared operator core). A single \texttt{OperatorCore} module ($K$ parallel read--transform--gate--write operators) supports three execution semantics: \textbf{scheduled} (external clock), \textbf{fixed-point} (deep equilibrium with contraction pressure), and \textbf{itinerant} (GLV competition dynamics producing heteroclinic channels). This screen tests all four variants on a compositional sequence-transform task with a common operator core, common task, and shared seed. The itinerant-fixed (hand-built chain) variant achieves near-perfect accuracy (99.6\%), dramatically outperforming scheduled (77.1\%) and fixed-point (56.0\%). The free-rho itinerant variant reaches 86.5\%, confirming that learned inhibition can produce useful itinerant dynamics without a hand-built chain.
\end{abstract}

\section{Summary}

\textbf{Experiment ID:} \texttt{20260720T000500Z-carom-gpu-screen-r1} \\
\textbf{Status:} completed (all four arms, exit code 0) \\
\textbf{Pod:} \texttt{00e42nckt4vm1v} (A100 SXM4 80GB, Secure Cloud) \\
\textbf{Cost bound:} USD 10.00 \\
\textbf{Approval:} Ben Goertzel, 2026-07-19

\section{Task}

The task is a compositional sequence-transform problem over a workspace of $n = 6$ slots in $\mathbb{Z}_8$. Programs consist of up to $L = 5$ commands drawn from a vocabulary of 10 tokens (8 distinct transforms plus 2 synonyms that force routing reuse). Commands include increment, double, negate, reverse, shift, swap, and cumulative-sum. A program is executed in presentation order; shorter programs are padded with \texttt{noop}.

The model receives the command sequence $C$ and the initial workspace $X$, and must predict the final workspace $Y = f_C(X)$.

\section{Architecture}

\subsection{Operator Core}

The shared \texttt{OperatorCore} implements $K = 16$ parallel operators. Each operator performs:

\begin{enumerate}
    \item \textbf{Read:} Compute attention-weighted read from workspace $w$ using concatenated content and positional features.
    \item \textbf{Transform:} Two-layer MLP with GELU activation.
    \item \textbf{Gate:} Sigmoid gate $g = \sigma(W_g z + b_g)$ controls write amplitude.
    \item \textbf{Write:} Produce update $u = g \cdot \tanh(z)$.
\end{enumerate}

The core parameters are:
\begin{itemize}
    \item $W_q, W_k, W_v \in \mathbb{R}^{K \times 3d \times d}$ (attention projections)
    \item $T_1 \in \mathbb{R}^{K \times 4d \times 2d}$, $T_2 \in \mathbb{R}^{K \times 2d \times d}$ (transform MLP)
    \item $W_g \in \mathbb{R}^{K \times d \times d}$, $b_g \in \mathbb{R}^{K \times d}$ (gate)
    \item LayerNorm on workspace content
\end{itemize}

Model dimension $d = 64$, $K = 16$ operators, total parameters $\approx 1.32\text{M}$.

\subsection{Execution Semantics}

\subsubsection*{Scheduled}
External clock iterates for $T = 5$ steps. At step $t$, routing weights $\beta_t = \text{softmax}(\text{route}(C_t) / \tau)$ select a mixture over operators:
\[
w \leftarrow w + \sum_{k=1}^{K} \beta_{t,k} \, u_k
\]

\subsubsection*{Fixed-Point (Deep Equilibrium)}
Starting from $w_0$, iterate the operator mixture for 8 sweeps with relaxation factor $\alpha = 0.5$:
\[
w \leftarrow (1 - \alpha) \, w + \alpha \left(w + \sum_k \beta_k \, u_k\right)
\]
Contraction pressure $\|r_{\text{last}}\|$ is added to the loss. Residuals $\|w_{t+1} - w_t\|$ are logged per sweep.

\subsubsection*{Itinerant (GLV Competition)}
Operators compete via Generalized Lotka--Volterra dynamics over $T_{\text{dyn}} = 70$ steps with step size $\Delta t = 0.1$, fatigue, and leak. The activity vector $a(t) \in \mathbb{R}^K$ evolves as:
\[
\dot{a}_k = a_k \left(r_k - \sum_j M_{kj} a_j - \phi_k\right)
\]
where $M$ is the inhibition matrix, $r_k$ is the intrinsic growth rate, and $\phi_k$ is a fatigue term. In the \textbf{fixed-chain} variant, $M$ is hand-built to produce a predetermined heteroclinic channel visiting operators in sequence. In the \textbf{free-rho} variant, $M$ is learned.

\section{Configuration}

\begin{table}[H]
\centering
\begin{tabular}{ll}
\toprule
\textbf{Parameter} & \textbf{Value} \\
\midrule
Seed & 7 \\
Steps & 4000 \\
Batch size & 128 \\
Model dimension $d$ & 64 \\
Operators $K$ & 16 \\
Optimizer & AdamW ($\text{lr}=2\text{e-}3$, $\text{wd}=1\text{e-}4$) \\
Scheduler & OneCycleLR \\
Gradient clip & 1.0 \\
Fixed-point sweeps & 8, $\alpha=0.5$ \\
Itinerant dynamics steps & 70, $\Delta t=0.1$ \\
Eval batch & 512 \\
Eval frequency & every 100 steps \\
Checkpoints & at steps 0, 1000, 2000, 3000, 3999 \\
GPU & NVIDIA A100-SXM4-80GB \\
\bottomrule
\end{tabular}
\end{table}

\section{Results}

\subsection{Final Accuracy}

\begin{table}[H]
\centering
\begin{tabular}{lccc}
\toprule
\textbf{Variant} & \textbf{Final Accuracy} & \textbf{Time (s)} & \textbf{Parameters} \\
\midrule
Scheduled & 0.771 & 82 & 1,322,472 \\
Fixed-Point & 0.560 & 203 & 1,322,792 \\
Itinerant-Fixed & 0.996 & 1,403 & 1,326,925 \\
Itinerant-Free & 0.865 & 1,417 & 1,326,925 \\
\bottomrule
\end{tabular}
\end{table}

\subsection{Learning Curves}

\subsubsection*{Scheduled}
\begin{lstlisting}
Step     0: loss=2.1033  acc=0.1416  |g|=0.3321  t=0s
Step   100: loss=2.0524  acc=0.1634  |g|=0.2158  t=3s
Step   500: loss=1.8065  acc=0.2751  |g|=0.3357  t=11s
Step  1000: loss=1.1893  acc=0.5417  |g|=0.7741  t=21s
Step  2000: loss=0.6308  acc=0.7803  |g|=0.2982  t=42s
Step  3000: loss=0.4709  acc=0.8167  |g|=0.1497  t=62s
Step  3999: loss=0.4092  acc=0.7965  |g|=0.2374  t=82s
\end{lstlisting}
Rapid early learning, plateaus around 0.78--0.82 after step 2000. Loss stabilizes near 0.41 but does not converge.

\subsubsection*{Fixed-Point}
\begin{lstlisting}
Step     0: loss=2.1040  acc=0.1416  |g|=0.3321  t=0s   res_last=0.0007
Step   100: loss=2.0598  acc=0.1576  |g|=0.2183  t=5s   res_last=0.0012
Step   500: loss=1.9347  acc=0.2161  |g|=0.2256  t=26s  res_last=0.0248
Step  1000: loss=1.7429  acc=0.3255  |g|=0.4034  t=51s  res_last=0.0413
Step  2000: loss=1.4228  acc=0.4492  |g|=0.3736  t=102s res_last=0.0515
Step  3000: loss=1.0729  acc=0.5999  |g|=0.4891  t=153s res_last=0.0542
Step  3999: loss=0.9462  acc=0.6081  |g|=0.4116  t=203s res_last=0.0541
\end{lstlisting}
Slowest learning. Convergence is poor despite contraction pressure. Residuals initially near zero, grow to $\sim$0.054 and stabilize---indicating the fixed-point iteration does not fully converge within 8 sweeps.

\textbf{Per-sweep residuals (final evaluation):}
\begin{lstlisting}
Sweep 0: 0.100161
Sweep 1: 0.089658
Sweep 2: 0.076235
Sweep 3: 0.066389
Sweep 4: 0.060115
Sweep 5: 0.055893
Sweep 6: 0.052934
Sweep 7: 0.050774
\end{lstlisting}
Residuals decrease monotonically (validating contraction), but the final residual (0.051) is non-negligible.

\subsubsection*{Itinerant-Fixed (Hand-Built Chain)}
\begin{lstlisting}
Step     0: loss=2.1032  acc=0.1416  |g|=0.3357  t=1s
Step   100: loss=2.0455  acc=0.1592  |g|=0.2204  t=35s
Step   500: loss=1.6875  acc=0.3141  |g|=0.5610  t=180s
Step  1000: loss=0.7018  acc=0.7214  |g|=0.8027  t=353s
Step  2000: loss=0.5029  acc=0.8288  |g|=0.6838  t=701s
Step  3000: loss=0.0333  acc=0.9935  |g|=0.7993  t=1055s
Step  3999: loss=0.0093  acc=0.9977  |g|=0.2981  t=1403s
\end{lstlisting}
Slow start (each step is $\sim$17$\times$ more expensive than scheduled due to 70-step GLV integration), but after step 2000, accuracy climbs steeply to near-perfect. Final loss 0.009 indicates near-convergence.

\textbf{Final activity mean (5 slots):} [0.0595, 0.0611, 0.0520, 0.0500, 0.0518] \\
\textbf{Active operators at final step:} 0 (all decayed)

\textbf{Itinerary order (sampled dynamics steps):}
\begin{lstlisting}
Step  0: active=[0]
Step  3: active=[0, 1]
Step  6: active=[1]
Step  9: active=[1, 2]
Step 12: active=[2]
Step 15: active=[2, 3]
Step 18: active=[3]
Step 21: active=[]
Steps 24-69: active=[] (all decayed)
\end{lstlisting}
Clean sequential visitation: operator $0 \to 0{+}1 \to 1 \to 1{+}2 \to 2 \to 2{+}3 \to 3 \to$ empty. This is the expected heteroclinic channel for the hand-built chain.

\subsubsection*{Itinerant-Free (Learned Inhibition)}
\begin{lstlisting}
Step     0: loss=2.1033  acc=0.1416  |g|=0.3355  t=1s
Step   100: loss=2.0478  acc=0.1644  |g|=0.2181  t=36s
Step   500: loss=1.8504  acc=0.2643  |g|=0.3268  t=175s
Step  1000: loss=1.3220  acc=0.4180  |g|=0.6413  t=353s
Step  2000: loss=0.5489  acc=0.8135  |g|=0.5365  t=707s
Step  3000: loss=0.2261  acc=0.8968  |g|=0.2709  t=1065s
Step  3999: loss=0.1943  acc=0.8802  |g|=0.1942  t=1417s
\end{lstlisting}
Learns faster than fixed-point but slower than itinerant-fixed. Plateaus around 0.88---good but clearly below the hand-built chain.

\textbf{Final activity mean (5 slots):} [0.0544, 0.0688, 0.0508, 0.0819, 0.9597] \\
\textbf{Active operators at final step:} 0.053 (near-zero)

The high activity in slot 4 (0.960) suggests one operator remains partially active---a different dynamics regime than the fixed-chain.

\textbf{Itinerary order (sampled dynamics steps):}
\begin{lstlisting}
Step  0: active=[0]
Step  3: active=[0, 1]
Step  6: active=[0, 1]
Step  9: active=[1]
Step 12: active=[1, 2]
Step 15: active=[1, 2]
Step 18: active=[2]
Step 21: active=[2]
Step 24: active=[2]
Step 27: active=[2, 3]
Step 30: active=[2, 3]
Step 33: active=[3]
Step 36: active=[3]
Step 39: active=[3]
Step 42: active=[]
Steps 45-69: active=[] (all decayed)
\end{lstlisting}
The learned inhibition produces a sequential visitation but with more overlap and dwell time at each operator. The sequence is correct but less crisp than the hand-built chain.

\section{Analysis}

\subsection{Itinerant Dynamics Are Decisive}
The itinerant-fixed variant's near-perfect accuracy (99.6\%) vs scheduled (77.1\%) demonstrates that sequential operator activation via heteroclinic channels is a powerful execution paradigm for compositional tasks. The key advantage: the itinerant dynamics enforce a temporal ordering on operator application that the scheduled variant must learn implicitly through routing alone.

\subsection{Fixed-Point Convergence Is the Bottleneck}
The fixed-point variant's poor performance (56.0\%) is explained by incomplete convergence: residuals decrease monotonically across sweeps but stabilize at $\sim$0.051, indicating the equilibrium is not reached within 8 sweeps. This is consistent with the CPU sandbox finding that residuals grow without contraction pressure. The contraction penalty is necessary but not sufficient.

\subsection{Learned Inhibition Works But Needs Refinement}
The free-rho variant (86.5\%) confirms that learned GLV inhibition can produce useful heteroclinic channels without a hand-built chain. The itinerary is sequentially correct but has more overlap and dwell time, suggesting the inhibition matrix needs stronger off-diagonal terms or longer dynamics steps for crisper transitions.

\subsection{Routing Bottleneck}
All variants share the same routing table $A$ mapping command tokens to operator mixtures. The two synonym tokens (\texttt{inc\_syn}, \texttt{rev\_syn}) force the same routing as their primary counterparts---a bottleneck that tests whether the operator core can generalize across shared routing. The itinerant-fixed variant's near-perfect accuracy suggests this bottleneck is not the limiting factor for compositional generalization.

\subsection{Compute Cost}
The itinerant variants are $\sim$17$\times$ more expensive per step than scheduled (35s/100 steps vs 2s/100 steps) due to 70-step GLV integration. This is the dominant cost. For larger tasks, amortizing the dynamics (e.g., learned step count or adaptive stopping) will be essential.

\section{Artifact Provenance}

\begin{table}[H]
\centering
\begin{tabular}{ll}
\toprule
\textbf{File} & \textbf{SHA-256} \\
\midrule
\texttt{summary.json} & \texttt{74b3e0796693fec3...0ff0b7f} \\
\texttt{scheduled.json} & \texttt{8dbc397adcf96f39...0e85315c} \\
\texttt{fixedpoint.json} & \texttt{fef5fc7d803d82f3...cd2c731e} \\
\texttt{itinerant-fixed.json} & \texttt{322d279ad80bc7e9...f8f2d4e52} \\
\texttt{itinerant-free.json} & \texttt{78d35512b66d54a3...035b321d1} \\
\texttt{mechanism\_screen.log} & \texttt{dd2fefac4dfdad73...68e0fa52} \\
\bottomrule
\end{tabular}
\end{table}

Artifacts stored at: \texttt{projects/carom/experiments/20260720T000500Z-carom-gpu-screen-r1/artifacts/carom\_screen/}

\section{Validity Gates}

\begin{itemize}
    \item[$\checkmark$] Fixed-point residuals decrease monotonically across sweeps (contraction confirmed)
    \item[$\checkmark$] Hand-built chain itinerancy is distinct from free-rho itinerancy (different activity patterns, different accuracy)
    \item[$\checkmark$] All four variants share common task, operator core, and seed
    \item[$\checkmark$] Exit code 0, all artifacts retrieved and verified
\end{itemize}

\section{Limitations}

\begin{enumerate}
    \item \textbf{Single seed.} No variance estimate. The ranking is clear but confidence intervals are unknown.
    \item \textbf{Toy task.} 6 slots, 8 values, 5 commands. Generalization to larger workspaces or longer programs is untested.
    \item \textbf{No implicit-gradient DEQ.} Fixed-point variant uses explicit sweeps, not backward-through-fixed-point. Performance may differ with implicit gradients.
    \item \textbf{No hyperparameter search.} All variants use the same lr, batch size, and schedule. Itinerant-specific tuning may improve free-rho.
    \item \textbf{One-off screen.} Not a broad claim of algorithmic superiority.
\end{enumerate}

\section{Conclusions}

\begin{table}[H]
\centering
\begin{tabular}{p{0.65\textwidth} p{0.3\textwidth}}
\toprule
\textbf{Finding} & \textbf{Status} \\
\midrule
Itinerant (fixed-chain) execution achieves near-perfect accuracy on compositional transforms & Observed \\
Scheduled execution plateaus at $\sim$77\% on the same task & Observed \\
Fixed-point execution is limited by incomplete convergence (residual $\sim$0.05) & Observed \\
Learned (free-rho) GLV inhibition produces a correct but less crisp itinerary & Observed \\
Temporal ordering via heteroclinic channels is the key mechanism for compositional generalization & Inferred \\
Stronger inhibition or adaptive dynamics could close the free-rho gap & Hypothesis \\
The routing bottleneck (synonym tokens) is not the limiting factor & Inferred \\
\bottomrule
\end{tabular}
\end{table}

\section{Next Steps}

\begin{enumerate}
    \item Multi-seed (3--5 seeds) confirmation of the accuracy ranking.
    \item Adaptive itinerant dynamics: learned step count or early stopping.
    \item Stronger free-rho inhibition: sweep fatigue/leak or initialize from fixed-chain prior.
    \item Larger tasks: longer programs, bigger workspace, more operators.
    \item Implicit-gradient DEQ training for the fixed-point variant.
    \item Scrambled-order discovery: can free-rho learn non-sequential itineraries?
\end{enumerate}

\vfill
\noindent\textit{Report generated 2026-07-20 04:50 PDT. Source: experiment r1 artifacts, CAROM repository \texttt{projects/carom/repos/carom/}.}

\end{document}
